I say that these latter are representable by open immersions. To see this, one applies criterion 3.1 (iii): $X_{i,j}$ satisfies the three exactness conditions (being a sheaf, commuting with $\varinjlim$ of rings and with adic $\varprojlim$ of artinian local rings), for $F,X_i,X_j$ satisfy them, and these conditions are stable under finite projective limits, in particular under fibered products; since $X_i\to F$ is a monomorphism, so is $v_{i,j}:X_{i,j}\to X_i$, which is deduced from it by the base change $X_j\to F$, and symmetrically $v_{j,i}$ is a monomorphism; finally the condition ``infinitesimally étale'' is also preserved by base change. This proves that one is under conditions 3.1 (iii).
% typo?: the printed base change X_j to F is retained.

We can now use the $X_i,X_{i,j},v_{i,j}$ and $w_{i,j}$ to construct in the usual way an $S$-prescheme $X$, such that the $X_i$ are identified with open subsets of $X$, the $X_{i,j}$ with the intersections $X_i\cap X_j$ and the $v_{i,j},w_{i,j}$ with the canonical immersions. Note that $X$ is also the quotient of
\[
\mathbf X=\coprod_i X_i
\]
by the equivalence relation $R=\coprod_{i,j}X_{i,j}$ (the two projections $v,w:R\rightrightarrows Y$ being defined by the $v_{i,j}$, resp. the $w_{i,j}$). More precisely, since $F$ is a sheaf for fpqc, the $u_i:X_i\to F$ come from a $u:\mathbf X\to F$, and $R$ is nothing other than the equivalence relation defined by $u$, $R=\mathbf X\times_F\mathbf X$; finally the quotient $X=Y/R$ is also a quotient in the category of sheaves for fpqc (and even in the category of sheaves for the Zariski topology): it suffices to use the definitions of ``quotient'' and ``sheaf'' to convince oneself of this. Consequently $u$ factors uniquely through a morphism
\[
u:X\to F,
\]
and this morphism is a monomorphism. It remains to show that it is an isomorphism. Since $F$ is of a local nature, one can suppose $S$ affine, and since moreover $F$ commutes with inductive limits of rings, it suffices to verify that for every $T$ affine of finite type over $S$, every morphism $T\to F$ factors through $X$ (\cf 3.4). For this, consider $G=X\times_F T\to T$; it suffices to prove that this is an isomorphism. Now, since $T$ is noetherian, one sees as above that it is an open immersion (N.B. $X\to F$ is infinitesimally étale, as follows at once from the fact that the induced morphisms $u_i:X_i\to F$ are). But by hypothesis $X\to F$ is set-theoretically surjective, and one sees immediately that this is a condition stable under base change, so $G\to T$ is set-theoretically surjective, hence an isomorphism since it is an open immersion. \hfill Q.E.D.
% typo?: the source alternates bold X and Y for the coproduct; kept.

\begin{proposition}{3.7}\label{XI.3.7}
Let $S$ be a locally noetherian prescheme, $I$ an increasingly filtered index set, $(T_i)_{i\in I}$ a projective system of $S$-preschemes locally of finite type, $T=\varprojlim T_i$ the projective limit functor, $F$ a functor over $S$, $u:F\to T$ an $S$-morphism. Suppose the following conditions hold:
\begin{enumerate}[label=\alph*)]
\item $F$ is a sheaf for the faithfully flat quasi-compact topology, commutes with inductive limits of rings, and with adic projective limits of artinian local rings.
\item The morphism $u:F\to T$ is a monomorphism.
\item[$\mathrm b'$)] The morphism $u:F\to T$ is infinitesimally étale.
\item For every point $\xi$ of $F$ with values in the spectrum of a field $k$, denoting by $\xi_i\in T_i(\Spec(k))$ its image and by $t_i$ the corresponding element of $T_i$, there exists an $i\in I$ such that for $j\ge i$ the transition morphism $p_{i,j}:T_j\to T_i$ is étale at $t_j$.
\item For every prescheme $X$ locally of finite type over $S$, and every $S$-morphism $X\to F$, the set of $x\in X$ at which this morphism is infinitesimally étale is open.
\end{enumerate}
Under these conditions, $F$ is representable by a prescheme locally of finite type over $S$.
\end{proposition}
Let us note at once that in the case which will concern us in the following \No, conditions c) and d) will be verified by means of the following corollary:
\begin{corollary}{3.8}\label{XI.3.8}
Assuming a), b), b'), conditions c) and d) are implied by the following:

c') The $T_i$ are smooth over $S$, and the transition morphisms $p_{i,j}:T_j\to T_i$ are smooth.

d') For every point $\xi$ of $F$ with values in the spectrum of a field $k$, let $t_i(\xi)$ be the element of $T_i$ defined by $\xi$, $d_i(\xi)$ the relative dimension of $T_i$ over $S$ at $t_i(\xi)$, and $d(\xi)=\operatorname{Sup}d_i(\xi)$. Then:
\begin{enumerate}[label=\arabic*$^\circ$)]
\item for every $\xi$ as above, one has $d(\xi)<+\infty$, and
\item for every prescheme $X$ locally of finite type over $S$, and every $S$-monomorphism $v:X\to F$, the function $x\mto d(\xi_x)$ on $X$ is locally constant (where, for $x\in X$, $\xi_x$ denotes the point of $F$ with values in $\kappa(x)$ induced by $v$).
\end{enumerate}
\end{corollary}

Let us prove 3.7. Place ourselves under the conditions of c), let $t$ be the element of $\operatorname{ens}(F)$ defined by $\xi$ (\cf 3.6) and put
\[
\OO_t=\varinjlim_i\OO_{T_i,t_i}.
\]
Then, using the condition stated in c), one sees easily that $\OO_t$ is a \emph{noetherian} local ring (EGA~$0_{\mathrm{IV}}$~10.3.1.3). Its residue field $\kappa(t)$ is the inductive limit of the residue fields $\kappa(t_i)$, and $k$ is an extension of it. If $\eta$ denotes the spectrum of $\kappa(t)$, $\xi$ that of $k$, one has a commutative diagram
\[
\xymatrix{\xi\ar[r]\ar[d]&F\ar[d]\\\eta\ar[r]&T,}
\]
whose definition is obvious and left to the reader. Since $\xi\to\eta$ is covering for the fpqc topology, $F$ is a sheaf for that topology by a), and $F\to T$ a monomorphism by b), one sees at once that the morphism $\eta\to T$ factors (uniquely) through a morphism $\eta\to F$. Let us now note the
\begin{lemma}{3.9}\label{XI.3.9}
Under the conditions of 3.7, let $T'=\Spec(A)$ be a noetherian local scheme, $T'_0=\Spec(k(A))$, $i:T'_0\to T'$ the canonical immersion, and suppose given a commutative square of morphisms
\[
\xymatrix{T'_0\ar[r]\ar[d]_i&F\ar[d]\\T'\ar[r]&T.}
\]
Then there exists a unique $T$-morphism $T'\to F$.
\end{lemma}
The proof is that of 3.1 (iii) $\Rightarrow$ (ii) (where the fact that the $S$ of the cited statement is representable was not used), using that $F$ is a sheaf for the fpqc topology, commutes with adic projective limits of artinian local rings (in this case the $A/\operatorname{rad}(A)^{n+1}$), and that $F\to T$ is a monomorphism and is infinitesimally étale.

We shall apply lemma 3.9 to the case where $T'=\Spec(\OO_t)$, hence $T'_0=\Spec(\kappa(t))$, noting that we have just constructed $T'_0\to F$, and that the canonical homomorphisms $\OO_{T_i,t_i}\to\OO_t$ define a projective system of morphisms $\Spec(\OO_t)\to T_i$, whence a canonical morphism $\Spec(\OO_t)\to T$. The commutativity of the corresponding square (x) is trivial (for, by definition of $T$, it suffices to verify it with $T$ replaced by $T_i$), whence a unique $T$-morphism
\[
v':T'=\Spec(\OO_t)\to F.
\]
Since $F$ commutes with inductive limits of rings, this morphism factors through
\[
v'_i:T'_i=\Spec(\OO_{T_i,t_i})\to F
\]
for large $i$. For such an $i$, one has
\[
T'\isomto T'_i\quad\ie\quad\OO_{T_i,t_i}\isomto\OO_t.
\]
Indeed, since $T'\to T'_i$ is faithfully flat quasi-compact and hence an effective epimorphism, it suffices to see that it is a monomorphism. Now if one has two morphisms $R\rightrightarrows T'$ (with $R$ representable) having the same composite with $T'\to T'_i$, they have the same composite with $T'\to T$ (which factors through $T'\to T'_i$), hence the same composite with $T'\to T_j$ for every $j$, hence with $T'\to T'_j$ for every $j$, and are therefore equal since $T'$ is the projective limit of the $T'_j$ in the category $\Sch$.

Consider the diagram
\[
\xymatrix{T'\ar[r]^{v'}\ar[d]&F\ar[d]\\T'_i\ar[r]\ar[ur]^{v'_i}&T_i,}
\]
where the unspecified arrows are the obvious arrows. The square is commutative and the upper triangle also, so since $T'\to T'_i$ is an epimorphism, it follows that the lower triangle is commutative. Now $\OO_{T_i,t_i}$ is the inductive limit of the affine rings of the affine open neighborhoods of $t_i$ in $T_i$, so since $F$ commutes with inductive limits of rings, $v'_i$ comes from a morphism
\[
v_t:U_t\to F
\]
from an open neighborhood $U_t$ of $t_i$ in $T_i$. One sees immediately that, on restricting this neighborhood if necessary, $v_t$ is necessarily a $T_i$-morphism ($T_i$ being locally of finite type over $S$, hence also commuting with inductive limits of rings). Then the composite of $v_t$ with $F\to T_i$ is a monomorphism, so $v_t$ is a monomorphism. I say that it is \emph{infinitesimally étale at $t_i$}. Since $F\to T$ is infinitesimally étale, it amounts to the same thing to say that the composite $T_i\to F\to T$ is infinitesimally étale at $t_i$, or again that the induced morphism $T'_i\to T$ is infinitesimally étale, or finally (since $T'\isomto T'_i$) that $T'\to T$ is infinitesimally étale, which is immediate, this morphism being the projective limit of the infinitesimally étale morphisms $T'_i\to T_i$.
% typo?: T_i to F to T, instead of U_t, kept as printed.

Let us now apply condition d) (which had not yet been used); it follows that $v_t$ is infinitesimally étale near $t$, so on replacing $U_t$ by a smaller open subset one can suppose that $v_t$ is an infinitesimally étale monomorphism.

For variable $t\in\operatorname{ens}(F)$, the family of morphisms $v_t$ falls under 3.5, which implies the conclusion of 3.7.

Let us now prove 3.8, supposing that conditions c') and d') of 3.8 hold. Then, with the notation of d'), one will have $d_i(\xi)=\text{constant}$ for large $i$, so the relative dimension of $T_j$ over $T_i$ at $t_j$, equal to $d_j(\xi)-d_i(\xi)$, is zero, hence $T_j\to T_i$ is étale at $t_j$, which proves condition c). Consequently the preceding proof applies to give, for every $t\in\operatorname{ens}(F)$, an index $i$, an open neighborhood $U_t$ of $t_i$ and a morphism $U_t\to F$ that is a monomorphism, infinitesimally étale at $t_i$, and everything amounts to proving that this morphism is infinitesimally étale near $t_i$. Now with the notation of d') $2^\circ$, where one takes $X=U_t$, one can suppose, on replacing $U_t$ by the connected component of $t_i$ in $U_t$, that for every $x\in U_t$ one has
\[
d(\xi_x)=d(\xi_{t_i})\quad\text{for every }x\in U_t.
\]
But $d(\xi_{t_i})=\text{relative dimension of }T_i\text{ over }S\text{ at }t_i$, and since the relative dimension of $T_i$ over $S$ remains constant on $U_t$, one will also have, for every $x\in U_t$:
\begin{equation}\tag{*}
d(\xi_x)=\text{relative dimension of }U_t\text{ over }S\text{ at }x.
\end{equation}
On the other hand, with the notation of the proof of 3.7, one sees at once that for every noetherian local prescheme $R'=\Spec(A)$, putting $R'_0=\Spec k(A)$, every morphism $R'\to F$ such that the induced morphism $R'_0\to F$ factors through $F$ (\ie taking the closed point of $R'$ to $t\in\operatorname{ens}(F)$) factors (obviously uniquely) through $T'$. (The proof is that of 3.1 or 3.9: one uses that $T'\to F$ is a monomorphism infinitesimally étale at $t$, and that $T'$ is a sheaf for fpqc commuting with adic $\varprojlim\ldots$). Applying this result to the morphism $\Spec(\OO_{U_t,x})=R'\to F$ induced by $v_t$ and to the point $t_x\in\operatorname{ens}(F)$ that is the image of the closed point of $R'$, \ie the image of $x$ under $v_t$, one finds a factorization
\[
\Spec(\OO_{U_t,x})\to\Spec(\OO_{t_x})\to F
\]
and since the second arrow is infinitesimally étale at $t_x$, to prove that the composite is so at $x$, it suffices to prove that the first is so at $t_x$.
% typo?: "factors through F" kept as printed.

Now, thanks to formula (x) above, it is a local $S$-homomorphism of the localizations, at two points $x,t_x$, of smooth $S$-preschemes $U_t,U_{t_x}$ of the same relative dimension $d$ over $S$ at $t,t_x$. This morphism is induced by an $S$-morphism
\[
w:U\to V
\]
where $U$ is an open neighborhood of $x$ in $U_t$, and where $V=U_{t_x}$ (EGA~I~6.5.1). Everything amounts to proving that this morphism is étale at $x$. Moreover $U$ and $V$ are endowed with monomorphisms $U\to F$, $V\to F$, and one sees at once that, on restricting $U$ further, $w$ is an $F$-morphism, so $w$ is a monomorphism. It now suffices to prove the
% typo?: the reference (x), while the displayed formula is (*), kept as printed.
\begin{lemma}{3.10}\label{XI.3.10}
Let $U,V$ be two smooth $S$-preschemes, of the same relative dimension $d$ over $S$, and $w:U\to V$ an $S$-morphism that is a monomorphism; then $w$ is an open immersion (and a fortiori is étale).
\end{lemma}
By SGA~I~5.7 one is reduced to the case where $S$ is the spectrum of a field, which one can suppose algebraically closed. By SGA~I~5.1 it suffices to prove that $w$ is étale, and it suffices to prove it at the closed points of $U$. Let $x$ be such a point, $y=w(x)$; then, taking a regular system of parameters $f_1,\ldots,f_d$ of $\OO_{V,y}$, one sees that $A=\OO_{U,x}/\sum f_i\OO_{U,x}$ is the trivial extension of $k=\OO_{V,y}/\sum f_i\OO_{V,y}$ (for, since $w$ is a monomorphism, so is the structural morphism $\Spec(A)\to\Spec(k)$ deduced from it by base change). Since $\OO_{U,x}$ is a regular local ring of dimension $d$, it follows that the $f_i$ form a regular system of parameters of this ring, which implies at once that $w$ is étale at $x$ and completes the proof of 3.10, hence of 3.8.

\begin{corollary}{3.11}\label{XI.3.11}
Under the conditions of 3.8, for every quasi-compact open subset $U$ of $F$ separated over $S$, there exists an $i\in I$ such that for every $j\ge i$ the morphism $u_j|U:U\to T_j$ is an open immersion. In particular, if the $T_i$ are quasi-affine over $S$, then every open subset $U$ of $F$ quasi-compact over $S$, \ie of finite type over $S$, is quasi-affine over $S$.
\end{corollary}
The proof of 3.7 shows that for every $t\in F$, there exists an $i\in I$ such that $u_i:F\to T_i$ is a local isomorphism at $t$, and then $u_j$ is a local isomorphism at $x$ for every $j\ge i$. By quasi-compactness, one can choose $i$ independent of $x\in U$. It remains to prove that for large $i$, $u_i|U:U\to T_i$ is a monomorphism. Now, since $U\to T$ is a monomorphism, one sees that the intersection of the equivalence relations $U\times_{T_i}U\subset U\times_S U$ is reduced to the diagonal, and since $U\times_S U$ is a noetherian prescheme and the $U\times_{T_i}U$ closed sub-preschemes, it follows that one of these $U\times_{T_i}U$ is already reduced to the diagonal, \ie $u_i|U$ is a monomorphism. This proves the first assertion in 3.11, and the second is an immediate consequence of it.

\begin{proposition}{3.12}\label{XI.3.12}
Let $S$ be a prescheme, $G$ an affine $S$-group prescheme.
\begin{enumerate}[label=\alph*)]
\item Let $F$ be the functor $\Sch^\circ_{/S}\to\Ens$ such that, for every $T$ over $S$,
\[
\begin{split}
F(T)=\text{set of subgroups of multiplicative type of }G_{T_S}\\
\text{that are finite over }T.
\end{split}
\]
Suppose $S$ locally noetherian or $G$ of finite presentation over $S$. Then the functor $F$ is representable and is affine over $S$. If $G$ is of finite presentation over $S$, then $F$ is locally of finite presentation over $S$.
% typo?: G_{T_S} kept as printed.
\item Let $H$ be an $S$-group prescheme of multiplicative type, and finite over $S$. Then $\uHom_{S\text{-gr}}(H,G)$ is representable. It is affine over $S$, and if $G$ is of finite type (resp. of finite presentation) over $S$, then so is $\uHom_{S\text{-gr}}(H,G)$.
\end{enumerate}
\end{proposition}
\begin{remark}{3.13}\label{XI.3.13}
Except for the precision that $\uHom_{S\text{-gr}}$ is affine, and in the case where $G$ is of finite presentation over $S$ (which will suffice for us), 3.12 is an immediate consequence of the theory of Hilbert Schemes (A.~Grothendieck, \textit{Techniques de construction et théorèmes d'existence en Géométrie Algébrique: IV Les Schémas de Hilbert}, Séminaire Bourbaki, May 1961, \No221). It even suffices that $G$ be quasi-projective over $S$; in case a), one can also represent the larger functor
\[
F'(T)=\left\{\begin{array}{l}
\text{set of group sub-preschemes of }G_T,\\
\text{flat, proper and of finite presentation over }T,
\end{array}\right\}
\]
(the canonical monomorphism $F\to F'$ is an ``open immersion'', as follows from criterion 3.1 and X~4.7 b), so the representability of $F'$ implies that $F$ is representable by an open subset of $F$); in case b), one can restrict oneself to supposing that $H$ is projective and of finite presentation over $S$. In both cases, one obtains a functor locally of finite presentation over $S$. In the present exposé, 3.12 is only a technical lemma for proving a key result in the following \No, so we shall sketch an easy direct proof of 3.12, not using Hilbert schemes.
% typo?: "an open subset of F" kept as printed.
\end{remark}

Let us first prove b). It will suffice that we prove that $\uHom_S(H,G)$ is representable (independently of any group structure on $H,G$), and has the additional properties stated for $\uHom_{S\text{-gr}}(H,G)$, for, using also the same result for $\uHom_S(H\times_S H,G)$, one describes the subfunctor $\uHom_{S\text{-gr}}(H,G)$ of $\uHom_S(H,G)$ by a finite projective limit (in fact, by means of fibered products) involving $G$, $\uHom_S(H,G)$ and $\uHom_S(H\times_S H,G)$, which we leave the reader to describe explicitly. On the other hand, one will have
\[
H=\Spec(\mathscr B),
\]
where $\mathscr B$ is a sheaf of algebras on $S$ that is locally free as a sheaf of modules (this is the only hypothesis on $H$ that we shall have to retain). Since the representability question considered is local on $S$, we can suppose $S$ affine with ring $A$. On the other hand, one will have $G=\Spec(C)$, where $C$ is an $A$-algebra. When $G=S[t]=G_0$, $t$ an indeterminate, the functor $\uHom$ is nothing other than
\[
T\mto\Gamma(T,\mathscr B_T),
\]
which is representable ($\mathscr B$ being locally free) by the vector bundle $\mathbf V(\mathscr B^\vee)$, where $\mathscr B^\vee$ is the sheaf of modules dual to $\mathscr B$. When $G=S[(t_i)]$, with $(t_i)$ a family (not necessarily finite) of indeterminates, one will therefore have $G=G_0^I$ (product over $S$ of a family of copies of $G_0$), which is representable by the affine scheme
\[
\uHom_S(H,G)^I=\mathbf V(\mathscr B^\vee)^I=\mathbf V(\mathscr B^{\vee(I)}).
\]
% typo?: Hom_S(H,G)^I kept as printed.
